%%% =======================================================================
%%% ISLS Extension Phase 5 v1.0.0
%%% C21: Adversarial Hypothesis Drill (PMHD)
%%% C22: Universal Artifact Intermediate Representation
%%% C23: Generative Synthesis Engine (Forge)
%%% Author: Sebastian Klemm
%%% Date: 16 March 2026
%%% Status: Implementation Specification for Claude Code
%%% =======================================================================
\documentclass[11pt,a4paper,twoside]{article}

\usepackage[utf8]{inputenc}
\usepackage{amsmath,amssymb,amsthm,mathtools}
\usepackage{algorithm,algorithmic}
\usepackage{booktabs,longtable,tabularx}
\usepackage{enumitem}
\usepackage{geometry}
\usepackage{hyperref}
\usepackage{cleveref}
\usepackage{xcolor}
\usepackage{fancyhdr}
\usepackage{listings}
\usepackage{tikz}
\usetikzlibrary{arrows.meta,positioning,shapes.geometric,fit}

\geometry{margin=2.5cm,top=3cm,bottom=3cm,headheight=14pt}

\theoremstyle{definition}
\newtheorem{definition}{Definition}[section]
\newtheorem{axiom}{Axiom}[section]
\newtheorem{requirement}{Requirement}[section]
\newtheorem{invariant}{Invariant}[section]

\theoremstyle{plain}
\newtheorem{theorem}{Theorem}[section]
\newtheorem{proposition}{Proposition}[section]
\newtheorem{lemma}{Lemma}[section]

\theoremstyle{remark}
\newtheorem{remark}{Remark}[section]

\newcommand{\ISLS}{\textsf{ISLS}}
\newcommand{\FORGE}{\textsf{FORGE}}
\newcommand{\PMHD}{\textsf{PMHD}}
\newcommand{\RR}{\mathbb{R}}
\newcommand{\NN}{\mathbb{N}}
\newcommand{\calH}{\mathcal{H}}
\newcommand{\calC}{\mathcal{C}}
\newcommand{\calQ}{\mathcal{Q}}
\newcommand{\calP}{\mathcal{P}}
\newcommand{\calM}{\mathcal{M}}
\newcommand{\calA}{\mathcal{A}}
\newcommand{\calE}{\mathcal{E}}
\newcommand{\calD}{\mathcal{D}}
\newcommand{\calF}{\mathcal{F}}
\newcommand{\calL}{\mathcal{L}}
\newcommand{\calS}{\mathcal{S}}
\DeclareMathOperator{\Dig}{Dig}
\DeclareMathOperator{\Canon}{Canon}
\DeclareMathOperator{\Drill}{Drill}
\DeclareMathOperator{\Oppose}{Oppose}
\DeclareMathOperator{\Commit}{Commit}
\DeclareMathOperator{\Synth}{Synth}
\DeclareMathOperator{\Eval}{Eval}
\newcommand{\norm}[1]{\left\lVert #1 \right\rVert}

\lstset{
  basicstyle=\ttfamily\small,
  breaklines=true,
  frame=single,
  backgroundcolor=\color{gray!5},
  rulecolor=\color{gray!30},
}

\pagestyle{fancy}
\fancyhf{}
\fancyhead[LE,RO]{\ISLS{} Extension Phase 5 v1.0.0}
\fancyhead[RE,LO]{Sebastian Klemm}
\fancyfoot[C]{\thepage}

\hypersetup{
  colorlinks=true,
  linkcolor=blue!60!black,
  pdftitle={ISLS Extension Phase 5 v1.0.0 — Generative Forge},
  pdfauthor={Sebastian Klemm},
}

\begin{document}

\begin{titlepage}
\centering
\vspace*{2cm}
{\Huge\bfseries \ISLS{} Extension Phase~5\\[0.3cm]
v1.0.0}\\[1.5cm]
{\Large C21: Adversarial Hypothesis Drill (\PMHD{})\\
C22: Universal Artifact IR\\
C23: Generative Synthesis Engine (\FORGE{})}\\[1cm]
{\large The Generative Counterpart:\\
From Intent to Evidence-Bound Artifact}\\[1cm]
{\normalsize Absorbing PMHD + ArtifactIR from MERKABA Factory\\
and Mining + Matrix from OmniCube}\\[2cm]
{\large Sebastian Klemm}\\[0.5cm]
{\large 16 March 2026}\\[1.5cm]

\begin{abstract}
\noindent
This document specifies the generative half of the \ISLS{} substrate:
three crates that transform structured intent into evidence-bound
artifacts through adversarial hypothesis drilling, universal
intermediate representation, and domain-specific synthesis.

C21 (\texttt{isls-pmhd}) implements the Polycentric Multi-Hypothesis
Drill: a generative process that produces hypotheses, tests them
adversarially against counterexamples, measures quality across six
axes (coherence, diversity, novelty, stability, robustness, coverage),
and commits survivors as Monoliths.

C22 (\texttt{isls-artifact-ir}) defines the universal, flat,
serializable intermediate representation that any Monolith is
transformed into before domain-specific synthesis.

C23 (\texttt{isls-forge}) is the synthesis engine that takes an
ArtifactIR, interprets it through a domain Matrix (Filament),
evaluates it against constraints, stores patterns in a persistent
memory, renders the final output, and validates the result through
\ISLS{}'s 8-gate cascade --- producing a Semantic Crystal that
contains the artifact with full evidence provenance from DecisionSpec
to final output.

Together with the analytical core (C1--C20), Phase~5 completes the
\ISLS{} substrate as a system that can both \emph{discover} programs
in data and \emph{generate} artifacts from intent, using identical
formal machinery for both directions.

\medskip\noindent\textbf{Status:} Implementation Specification.\\
\textbf{Parent:} ISLS v1.0.0 + Extensions Phase 1--4.\\
\textbf{Absorbs:} MERKABA \texttt{domain-artifact} (PMHD, ArtifactIR,
  Synth, Eval, Memory, Render), OmniCube \texttt{blueprint-forge}
  (Mining, Matrix, Materializer).\\
\textbf{Depends on:} C12 (Registry), C13 (Manifest), C14 (Capsule),
  C17 (Store), C19 (Gateway).
\end{abstract}
\end{titlepage}

\tableofcontents
\newpage


%%% =====================================================================
\part{Foundations}\label{part:foundations}

\section{The Two Directions}\label{sec:twodirections}

\ISLS{} until Phase~4 is purely analytical: it observes data, discovers
patterns, and crystallizes them. Phase~5 adds the generative direction:
it accepts structured intent, drills it adversarially, synthesizes an
artifact, and crystallizes the result through the same 8-gate cascade.

\begin{definition}[Analytical Direction]\label{def:analytical}
\[
  \text{World} \xrightarrow{\text{observe}} G
  \xrightarrow{\text{extract}} P
  \xrightarrow{\text{validate}} C_{\text{discovery}}
\]
Input: observations. Output: crystal encoding a discovered program.
\end{definition}

\begin{definition}[Generative Direction]\label{def:generative}
\[
  \text{Intent} \xrightarrow{\text{drill}} M
  \xrightarrow{\text{IR}} A
  \xrightarrow{\text{synth}} O
  \xrightarrow{\text{validate}} C_{\text{forge}}
\]
Input: DecisionSpec. Output: crystal encoding a generated artifact.
\end{definition}

\begin{axiom}[Substrate Identity]\label{ax:substrate}
Both directions use the same state space ($\RR^5$), the same
canonicalization (JCS + SHA-256), the same Run Descriptor, the same
Evidence Chain, the same Crystal format, and the same 8-gate cascade.
A crystal produced by the forge is structurally indistinguishable from
a crystal produced by discovery. Only the \texttt{program\_id} field
in the manifest distinguishes provenance.
\end{axiom}

\begin{axiom}[Adversarial Quality]\label{ax:adversarial}
No artifact may be committed without surviving adversarial opposition.
The PMHD drill is not a generator with a quality filter --- it is a
dialectical process where hypotheses are tested against structured
counterevidence before commitment.
\end{axiom}

\begin{axiom}[Memory Persistence]\label{ax:memory}
The pattern memory is append-only and persistent across runs. Every
committed monolith enriches the memory for future runs. The memory
is stored in the \ISLS{} archive (C7) and store (C17), not in an
ephemeral in-process structure.
\end{axiom}


\section{Absorption Map}\label{sec:absorb}

\begin{tabularx}{\textwidth}{l l X}
\toprule
\textbf{Source} & \textbf{Target} & \textbf{Transformation} \\
\midrule
MERKABA \texttt{pmhd/} & C21 \texttt{isls-pmhd} & DrillEngine,
  Hypothesis, Counterexample, QualityMetrics, PmhdMonolith,
  Provenance. Adapted to use \ISLS{} types, registry, and PoR
  from C5. \\
MERKABA \texttt{artifact\_ir} & C22 \texttt{isls-artifact-ir} &
  ArtifactIR, ArtifactHeader, Component, Interface, Constraint,
  Metrics, Provenance, Delta. Unchanged structurally; retyped to
  use Hash256 and \ISLS{} serialization. \\
MERKABA \texttt{synth/eval/memory/render} & C23 \texttt{isls-forge} &
  Synthesizers, Evaluators, PatternLibrary, Emitters. Extended with
  8-gate validation and crystal output. \\
OmniCube \texttt{qops} & C21 mining strategies & Greedy, Stochastic,
  Beam, Evolutionary search strategies for hypothesis seed
  generation. \\
OmniCube \texttt{blueprint-forge} & C23 domain matrices & Matrix
  concept: domain-specific interpretation of abstract signatures
  into concrete artifacts. \\
\bottomrule
\end{tabularx}


%%% =====================================================================
\part{Crate C21: \texttt{isls-pmhd}}\label{part:pmhd}

\section{Purpose}\label{sec:pmhd:purpose}

The PMHD crate implements the adversarial hypothesis drill: a
deterministic process that generates, opposes, evaluates, and commits
structured hypotheses about what an artifact should be.


\section{Core Objects}\label{sec:pmhd:objects}

\begin{definition}[DecisionSpec]\label{def:decisionspec}
A DecisionSpec $\calD$ is the input to the generative direction:
\[
  \calD = \langle \mathtt{id}, \mathtt{intent}, \mathtt{goals},
          \mathtt{constraints}, \mathtt{seeds},
          \mathtt{domain}, \mathtt{config} \rangle
\]
where:
\begin{itemize}[nosep]
  \item \texttt{id}: content-address of the canonical spec,
  \item \texttt{intent}: natural-language or structured description of
        the desired artifact,
  \item \texttt{goals}: key-value map of target quality metrics
        (e.g., \texttt{\{coherence: 0.7, robustness: 0.8\}}),
  \item \texttt{constraints}: hard constraints that must be satisfied
        (list of predicate strings),
  \item \texttt{seeds}: optional seed crystals from the \ISLS{} archive
        to guide hypothesis generation,
  \item \texttt{domain}: target domain identifier
        (e.g., \texttt{"rust"}, \texttt{"api"}, \texttt{"workflow"}),
  \item \texttt{config}: PMHD configuration (tick count, pool size,
        opposition strength).
\end{itemize}
\end{definition}

\begin{definition}[Hypothesis]\label{def:hypothesis}
A hypothesis $h \in \calH$ is a structured claim:
\[
  h = \langle \mathtt{id}, \mathtt{claim}, \mathtt{assumptions},
      \mathtt{evidence}, \mathtt{generation} \rangle
\]
where:
\begin{itemize}[nosep]
  \item \texttt{id} = SHA-256(claim $\|$ sorted(assumptions)),
  \item \texttt{claim}: the core assertion (what the artifact should be),
  \item \texttt{assumptions}: ordered list of supporting assumptions,
  \item \texttt{evidence}: list of evidence items supporting the claim,
  \item \texttt{generation}: the drill tick at which this hypothesis
        was created.
\end{itemize}
\end{definition}

\begin{definition}[Counterexample]\label{def:counterexample}
A counterexample $c \in \calC$ is a structured objection:
\[
  c = \langle \mathtt{id}, \mathtt{target\_id}, \mathtt{reason},
      \mathtt{severity} \rangle
\]
where $\mathtt{severity} \in [0, 255]$ quantifies the strength of the
objection.
\end{definition}

\begin{definition}[Quality Metrics]\label{def:quality}
The quality of a hypothesis is measured across six axes:
\[
  \calQ(h) = \langle q_{\mathrm{coh}}, q_{\mathrm{div}},
  q_{\mathrm{nov}}, q_{\mathrm{stab}}, q_{\mathrm{rob}},
  q_{\mathrm{cov}} \rangle \in [0,1]^6
\]
\begin{itemize}[nosep]
  \item $q_{\mathrm{coh}}$: \textbf{Coherence} --- internal consistency
        of the hypothesis and its evidence.
  \item $q_{\mathrm{div}}$: \textbf{Diversity} --- breadth of
        perspectives explored in the drill pool.
  \item $q_{\mathrm{nov}}$: \textbf{Novelty} --- distance from
        previously committed monoliths in the pattern memory.
  \item $q_{\mathrm{stab}}$: \textbf{Stability} --- resilience under
        repeated adversarial opposition across ticks.
  \item $q_{\mathrm{rob}}$: \textbf{Robustness} --- survival rate
        against counterexamples weighted by severity.
  \item $q_{\mathrm{cov}}$: \textbf{Coverage} --- fraction of the
        goal space addressed by the hypothesis.
\end{itemize}
\end{definition}

\begin{definition}[Monolith]\label{def:monolith}
A PmhdMonolith $m \in \calM$ is the committed output of a drill cycle
that survived all quality thresholds and the PoR gate:
\[
  m = \langle \mathtt{id}, h^*, \mathtt{excalibration},
      \{c_1, \ldots, c_k\}, \calQ(h^*), \mathtt{provenance} \rangle
\]
where $h^*$ is the surviving core hypothesis, \texttt{excalibration}
is the adversarial sharpening vector, and provenance contains seed,
config hash, tick range, active set stats, and PoR evidence.
\end{definition}


\section{Drill Algorithm}\label{sec:pmhd:algo}

\begin{algorithm}[H]
\caption{\textsc{PMHD\_Drill}: Adversarial hypothesis drill}
\label{alg:drill}
\begin{algorithmic}[1]
\REQUIRE DecisionSpec $\calD$, Config $\theta$, PatternMemory $\calP$
\ENSURE List of PmhdMonoliths $\{m_1, \ldots, m_k\}$
\STATE $\mathrm{pool} \gets \textsc{SeedHypotheses}(\calD, \calP, \theta)$
\FOR{$t = 1$ \TO $\theta.\mathrm{ticks}$}
  \STATE \textit{// Phase 1: Generate}
  \STATE $\mathrm{pool} \gets \mathrm{pool} \cup
    \textsc{Mutate}(\mathrm{pool}, \calD, t)$
  \STATE \textit{// Phase 2: Oppose}
  \FOR{each $h \in \mathrm{pool}$}
    \STATE $\mathrm{counters} \gets \textsc{GenerateCounterexamples}(h, \calD, t)$
    \STATE $h.\mathrm{evidence} \gets h.\mathrm{evidence} \cup
      \textsc{Defend}(h, \mathrm{counters})$
  \ENDFOR
  \STATE \textit{// Phase 3: Evaluate}
  \FOR{each $h \in \mathrm{pool}$}
    \STATE $\calQ(h) \gets \textsc{MeasureQuality}(h, \mathrm{pool}, \calP)$
  \ENDFOR
  \STATE \textit{// Phase 4: Select and Prune}
  \STATE $\mathrm{pool} \gets \textsc{TournamentSelect}(\mathrm{pool},
    \theta.\mathrm{pool\_size})$
  \STATE \textit{// Phase 5: Commit Gate}
  \FOR{each $h \in \mathrm{pool}$ where $\calQ(h) \ge \theta.\mathrm{thresholds}$}
    \IF{$\textsc{PoR\_Gate}(h, \calQ(h), \calD) = \mathrm{FIRE}$}
      \STATE $m \gets \textsc{CommitMonolith}(h, \calQ(h), t, \calD)$
      \STATE $\mathrm{monoliths} \gets \mathrm{monoliths} \cup \{m\}$
      \STATE $\calP.\mathrm{add}(m)$
    \ENDIF
  \ENDFOR
\ENDFOR
\RETURN $\mathrm{monoliths}$
\end{algorithmic}
\end{algorithm}

\begin{definition}[Seed Strategies]\label{def:seedstrat}
Hypothesis seeding uses strategies absorbed from OmniCube QOPS:
\begin{enumerate}[nosep,label=(\roman*)]
  \item \textbf{Greedy}: seed from the highest-resonance crystals in
        the pattern memory that match the DecisionSpec domain.
  \item \textbf{Stochastic}: deterministic pseudo-random generation
        from the RD seed, sampling the 5D signature space.
  \item \textbf{Beam}: maintain $k$ best hypotheses and expand each.
  \item \textbf{Evolutionary}: mutation + crossover on the hypothesis
        pool with tournament selection.
  \item \textbf{Hybrid}: combine greedy seeding with evolutionary
        expansion (default).
\end{enumerate}
All strategies are deterministic under fixed $(\calD, \mathrm{RD})$.
\end{definition}

\begin{definition}[Commit Gate]\label{def:commitgate}
The PMHD commit gate fires if and only if:
\[
  G_{\PMHD}(h) = \mathbf{1}\!\left[
    q_{\mathrm{coh}} \ge \theta_{\mathrm{coh}} \;\wedge\;
    q_{\mathrm{div}} \ge \theta_{\mathrm{div}} \;\wedge\;
    q_{\mathrm{nov}} \ge \theta_{\mathrm{nov}} \;\wedge\;
    q_{\mathrm{stab}} \ge \theta_{\mathrm{stab}} \;\wedge\;
    q_{\mathrm{rob}} \ge \theta_{\mathrm{rob}} \;\wedge\;
    q_{\mathrm{cov}} \ge \theta_{\mathrm{cov}} \;\wedge\;
    \mathrm{PoR}(h)
  \right]
\]
The PoR check uses the existing \ISLS{} PoR FSM (C5) applied to the
drill state trajectory.
\end{definition}


\section{Configuration}\label{sec:pmhd:config}

\begin{lstlisting}[title={PMHD Configuration}]
pmhd:
  ticks: 5000
  pool_size: 50
  opposition_strength: 0.7   # severity multiplier
  mutation_rate: 0.3
  crossover_rate: 0.2
  seed_strategy: "hybrid"
  thresholds:
    coherence: 0.6
    diversity: 0.4
    novelty: 0.3
    stability: 0.5
    robustness: 0.6
    coverage: 0.4
  commit_budget: 20           # max monoliths per run
  use_pattern_memory: true
\end{lstlisting}


\section{Rust API}\label{sec:pmhd:api}

\begin{lstlisting}
// isls-pmhd/src/lib.rs

pub struct DecisionSpec {
    pub id: Hash256,
    pub intent: String,
    pub goals: BTreeMap<String, f64>,
    pub constraints: Vec<String>,
    pub seeds: Vec<Hash256>,        // crystal IDs from archive
    pub domain: String,
    pub config: PmhdConfig,
}

pub struct Hypothesis {
    pub id: String,                 // SHA-256(claim || assumptions)
    pub claim: String,
    pub assumptions: Vec<String>,
    pub evidence: Vec<Evidence>,
    pub generation: u64,
}

pub struct Counterexample {
    pub id: String,
    pub target_id: String,
    pub reason: String,
    pub severity: u8,
}

pub struct QualityMetrics {
    pub coherence: f64,
    pub diversity: f64,
    pub novelty: f64,
    pub stability: f64,
    pub robustness: f64,
    pub coverage: f64,
}

pub struct PmhdMonolith {
    pub id: String,
    pub core_hypothesis: Hypothesis,
    pub excalibration_vector: Vec<f64>,
    pub counterexamples: Vec<Counterexample>,
    pub quality: QualityMetrics,
    pub provenance: MonolithProvenance,
}

pub struct DrillEngine {
    config: PmhdConfig,
    pool: Vec<Hypothesis>,
    monoliths: Vec<PmhdMonolith>,
    pattern_memory: PatternMemory,
    tick: u64,
}

pub enum SeedStrategy {
    Greedy,
    Stochastic,
    Beam { width: usize },
    Evolutionary,
    Hybrid,
}

impl DrillEngine {
    pub fn new(config: PmhdConfig) -> Self;
    pub fn with_memory(config: PmhdConfig, memory: PatternMemory) -> Self;
    pub fn drill(&mut self, spec: &DecisionSpec) -> DrillResult;
    pub fn step(&mut self, spec: &DecisionSpec) -> StepResult;
    pub fn monoliths(&self) -> &[PmhdMonolith];
    pub fn pattern_memory(&self) -> &PatternMemory;
}

pub struct DrillResult {
    pub monoliths: Vec<PmhdMonolith>,
    pub ticks_executed: u64,
    pub pool_final_size: usize,
    pub quality_history: Vec<Vec<QualityMetrics>>,
    pub commit_count: usize,
}
\end{lstlisting}


\section{Acceptance Tests}\label{sec:pmhd:tests}

\begin{longtable}{l l p{6.5cm}}
\toprule
\textbf{ID} & \textbf{Name} & \textbf{Procedure} \\
\midrule
\endfirsthead \midrule \endhead \bottomrule \endfoot
AT-P1 & Drill determinism & Run drill twice with same spec and RD;
  verify identical monolith IDs. \\
AT-P2 & Opposition effect & Run with opposition\_strength=0 vs.
  opposition\_strength=0.9; verify robustness scores differ. \\
AT-P3 & Commit gate & Set all thresholds to 1.0; verify zero
  monoliths committed. Set to 0.0; verify all committed. \\
AT-P4 & Quality metrics range & Verify all 6 metrics in $[0,1]$ for
  every hypothesis at every tick. \\
AT-P5 & Pattern memory growth & Commit 3 monoliths; verify memory
  contains 3 patterns. Run again with same spec; verify novelty scores
  reflect existing patterns. \\
AT-P6 & Seed strategies & Run with each strategy; verify all produce
  at least 1 monolith. Verify determinism for each. \\
AT-P7 & Hypothesis ID determinism & Create hypothesis with same claim
  and assumptions; verify identical ID. \\
AT-P8 & Provenance completeness & Verify monolith provenance contains
  seed, config\_hash, tick range, PoR evidence. \\
\end{longtable}


%%% =====================================================================
\part{Crate C22: \texttt{isls-artifact-ir}}\label{part:ir}

\section{Purpose}\label{sec:ir:purpose}

The ArtifactIR is the universal, flat, serializable intermediate
representation that bridges the gap between the abstract monolith
(the adversarially-tested hypothesis) and the concrete domain-specific
output. It is domain-agnostic: the same IR structure represents a Rust
module, an API specification, a workflow definition, a database schema,
or any other structured artifact.


\section{Data Model}\label{sec:ir:model}

\begin{definition}[ArtifactIR]\label{def:artifactir}
An ArtifactIR $\calA$ is a flat structure:
\[
  \calA = \langle \mathtt{header}, \mathtt{components},
          \mathtt{interfaces}, \mathtt{constraints},
          \mathtt{metrics}, \mathtt{provenance},
          \mathtt{deltas}, \mathtt{extra} \rangle
\]
\end{definition}

\begin{definition}[ArtifactHeader]\label{def:artheader}
\begin{lstlisting}
pub struct ArtifactHeader {
    pub artifact_id: Hash256,      // content-addressed
    pub version: String,           // semantic version
    pub timestamp_tick: u64,       // drill tick of commitment
    pub layer_index: u32,          // hierarchy depth
    pub por_decision: Option<String>, // PoR evidence JSON
    pub source_monolith_id: String,   // link to PmhdMonolith
    pub domain: String,               // target domain
}
\end{lstlisting}
\end{definition}

\begin{definition}[Component]\label{def:component}
A component is an atomic unit of the artifact:
\begin{lstlisting}
pub struct Component {
    pub id: String,
    pub kind: String,       // "function", "type", "endpoint", etc.
    pub name: String,       // human-readable name
    pub content: String,    // canonical JSON of the component body
    pub dependencies: Vec<String>,  // component IDs this depends on
    pub signature: FiveDState,      // 5D embedding of this component
}
\end{lstlisting}
\end{definition}

\begin{definition}[Interface]\label{def:interface}
An interface defines a contract between components:
\begin{lstlisting}
pub struct Interface {
    pub id: String,
    pub provider: String,    // component ID that provides
    pub consumer: String,    // component ID that consumes
    pub contract: String,    // canonical JSON of the contract
    pub direction: Direction,
}

pub enum Direction { Unidirectional, Bidirectional }
\end{lstlisting}
\end{definition}

\begin{definition}[ArtifactMetrics]\label{def:artmetrics}
Quality metrics propagated from the monolith plus synthesis metrics:
\begin{lstlisting}
pub struct ArtifactMetrics {
    // From QualityMetrics (monolith)
    pub coherence: f64,
    pub diversity: f64,
    pub novelty: f64,
    pub stability: f64,
    pub robustness: f64,
    pub coverage: f64,
    // Synthesis-specific
    pub component_count: usize,
    pub interface_count: usize,
    pub constraint_satisfaction: f64,  // fraction of constraints met
    pub quality_score: f64,            // composite score
}
\end{lstlisting}
\end{definition}

\begin{definition}[ArtifactProvenance]\label{def:artprov}
Full traceability from DecisionSpec to artifact:
\begin{lstlisting}
pub struct ArtifactProvenance {
    pub decision_spec_id: Hash256,
    pub monolith_id: String,
    pub seed: u64,
    pub config_hash: String,
    pub tick_range: [u64; 2],
    pub drill_strategy: String,
    pub por_evidence: Option<String>,
    pub pattern_memory_size: usize,
}
\end{lstlisting}
\end{definition}

\begin{definition}[Build from Monolith]\label{def:buildir}
The IR is constructed deterministically from a monolith:
\begin{lstlisting}
impl ArtifactIR {
    pub fn build_from_monolith(
        monolith: &PmhdMonolith,
        spec: &DecisionSpec,
        layer: u32,
    ) -> Result<ArtifactIR>;
}
\end{lstlisting}
The construction is deterministic: same monolith and spec produce
identical ArtifactIR (content-addressed).
\end{definition}


\section{Acceptance Tests}\label{sec:ir:tests}

\begin{longtable}{l l p{6.5cm}}
\toprule
\textbf{ID} & \textbf{Name} & \textbf{Procedure} \\
\midrule
\endfirsthead \midrule \endhead \bottomrule \endfoot
AT-IR1 & IR determinism & Build IR from same monolith twice; verify
  identical artifact\_id. \\
AT-IR2 & Serde round-trip & Serialize IR to JSON; deserialize; verify
  equal. \\
AT-IR3 & Provenance link & Build IR; verify provenance.monolith\_id
  matches source monolith. \\
AT-IR4 & Component signature & Build IR; verify every component has a
  non-zero FiveDState signature. \\
\end{longtable}


%%% =====================================================================
\part{Crate C23: \texttt{isls-forge}}\label{part:forge}

\section{Purpose}\label{sec:forge:purpose}

The Forge is the synthesis engine that takes an ArtifactIR and produces
a concrete, domain-specific output --- validated through \ISLS{}'s
8-gate cascade and archived as a Semantic Crystal.


\section{Pipeline}\label{sec:forge:pipeline}

\begin{lstlisting}[title={Forge Pipeline}]
DecisionSpec
  |
  v
[C21: PMHD Drill] --> PmhdMonolith
  |
  v
[C22: ArtifactIR] --> ArtifactIR
  |
  v
[C23: Matrix Interpret] --> Domain-Specific IR
  |
  v
[C23: Synthesize] --> Raw Output (code, spec, schema...)
  |
  v
[C23: Evaluate] --> Constraint Check + Quality Bounds
  |
  v
[C23: Pattern Memory] --> Store for future runs
  |
  v
[ISLS: 8-Gate Validate] --> Formal validation
  |
  v
[ISLS: Crystal Commit] --> Semantic Crystal containing the artifact
  |
  v
[C23: Emit] --> Domain-specific files on disk / via Gateway
\end{lstlisting}


\section{Domain Matrices (Filaments)}\label{sec:forge:matrix}

\begin{definition}[Matrix]\label{def:matrix}
A Matrix $\calF$ is a domain-specific interpreter that maps
ArtifactIR components to concrete output structures:
\begin{lstlisting}
pub trait Matrix: Send + Sync {
    fn name(&self) -> &str;
    fn domain(&self) -> &str;
    fn interpret(
        &self,
        ir: &ArtifactIR,
        config: &MatrixConfig,
    ) -> Result<SynthesisInput>;
    fn supported_kinds(&self) -> Vec<String>;
}
\end{lstlisting}
\end{definition}

\begin{definition}[Built-in Matrices]\label{def:builtinmatrix}
Phase~5 ships with four built-in matrices:

\begin{tabularx}{\textwidth}{l l X}
\toprule
\textbf{Matrix} & \textbf{Domain} & \textbf{Output} \\
\midrule
RustModuleMatrix & \texttt{rust} & Rust module skeletons with types,
  traits, impls, tests. Components map to functions/types, interfaces
  to trait bounds. \\
HttpApiMatrix & \texttt{api} & OpenAPI 3.0 specifications. Components
  map to endpoints, interfaces to request/response schemas. \\
WorkflowMatrix & \texttt{workflow} & YAML workflow definitions
  (GitHub Actions / generic). Components map to steps, interfaces to
  data flows between steps. \\
SchemaMatrix & \texttt{schema} & JSON Schema definitions. Components
  map to type definitions, interfaces to \texttt{\$ref} relations. \\
\bottomrule
\end{tabularx}

Custom matrices are registered at runtime via the Matrix trait.
\end{definition}


\section{Synthesis}\label{sec:forge:synth}

\begin{definition}[Synthesizer]\label{def:synth}
A synthesizer takes the matrix-interpreted structure and produces
the raw output string:
\begin{lstlisting}
pub trait Synthesizer: Send + Sync {
    fn synthesize(
        &self,
        input: &SynthesisInput,
        config: &ForgeConfig,
    ) -> Result<SynthesisOutput>;
}

pub struct SynthesisOutput {
    pub content: String,           // the generated artifact text
    pub format: String,            // "rust", "yaml", "json", etc.
    pub components_rendered: usize,
    pub triplet: PerformanceTriplet,
}

pub struct PerformanceTriplet {
    pub psi: f64,    // precision
    pub rho: f64,    // recall / completeness
    pub omega: f64,  // structural quality
}
\end{lstlisting}
\end{definition}


\section{Evaluation}\label{sec:forge:eval}

\begin{definition}[Evaluator Stack]\label{def:evalstack}
The evaluation stack checks the synthesized output against multiple
criteria:
\begin{enumerate}[nosep,label=(\roman*)]
  \item \textbf{Constraint Evaluator}: verifies all hard constraints
        from the DecisionSpec are satisfied.
  \item \textbf{Quality Bounds Evaluator}: verifies the
        PerformanceTriplet meets minimum thresholds.
  \item \textbf{Homeostasis Bridge}: checks that the synthesis did not
        degrade quality metrics below the monolith's original scores.
\end{enumerate}
\begin{lstlisting}
pub struct EvaluationResult {
    pub constraints_met: usize,
    pub constraints_total: usize,
    pub quality_passed: bool,
    pub homeostasis_passed: bool,
    pub details: Vec<EvalDetail>,
}
\end{lstlisting}
\end{definition}


\section{Pattern Memory}\label{sec:forge:memory}

\begin{definition}[PatternMemory]\label{def:patternmem}
The pattern memory is a persistent, append-only library of committed
monolith patterns:
\begin{lstlisting}
pub struct PatternEntry {
    pub monolith_id: String,
    pub domain: String,
    pub quality: QualityMetrics,
    pub signature: FiveDState,   // 5D embedding of the pattern
    pub component_kinds: Vec<String>,
    pub timestamp: f64,
}

pub struct PatternMemory {
    entries: Vec<PatternEntry>,  // loaded from store at startup
}

impl PatternMemory {
    pub fn load(store: &IslandStore) -> Result<Self>;
    pub fn add(&mut self, entry: PatternEntry, store: &IslandStore) -> Result<()>;
    pub fn find_similar(
        &self, signature: &FiveDState, k: usize
    ) -> Vec<&PatternEntry>;
    pub fn novelty_score(&self, candidate: &FiveDState) -> f64;
    pub fn len(&self) -> usize;
}
\end{lstlisting}
The memory is persisted in \texttt{isls-store} (C17) as a dedicated
table. Every \texttt{add} writes to both the in-memory vector and the
database.
\end{definition}


\section{8-Gate Crystal Validation}\label{sec:forge:validate}

\begin{requirement}[Forge Crystal Validation]\label{req:forgeval}
Before the forge emits an artifact, the synthesized output
\textbf{MUST} be packaged as a Semantic Crystal and validated through
the standard 8-gate cascade (content\_address, dual\_consensus,
evidence\_chain, free\_energy, gate\_kairos, immutability,
operator\_versions, por\_trace). The crystal's
\texttt{constraint\_program} field contains the ArtifactIR serialized
as a canonical JSON blob. The crystal's \texttt{topology\_signature}
encodes the component dependency graph.
\end{requirement}

\begin{theorem}[Forge-Discovery Symmetry]\label{thm:forgesym}
A crystal produced by the forge and a crystal produced by analytical
discovery are both valid \ISLS{} crystals that pass the same 8 formal
checks. They differ only in provenance: the forge crystal's manifest
has \texttt{program\_id = "forge:<spec\_id>"} while the discovery
crystal has \texttt{program\_id = "discovery"}.
\end{theorem}

\begin{proof}
By construction: the forge crystal is built using the same
\texttt{build\_crystal\_with\_id} function, the same evidence chain
builder, the same PoR FSM, and the same 8-check validation harness.
The crystal format is identical. Only the program\_id metadata
differs.
\end{proof}


\section{Emitters}\label{sec:forge:emit}

\begin{definition}[Emitter]\label{def:emitter}
An emitter writes the final artifact to disk, gateway, or external
system:
\begin{lstlisting}
pub trait Emitter: Send + Sync {
    fn name(&self) -> &str;
    fn emit(
        &self,
        output: &SynthesisOutput,
        crystal: &SemanticCrystal,
        config: &EmitConfig,
    ) -> Result<EmitResult>;
}

pub struct EmitResult {
    pub path: Option<String>,
    pub crystal_id: Hash256,
    pub bytes_written: usize,
}
\end{lstlisting}
Built-in emitters: FileEmitter (write to disk), GatewayEmitter (POST
to gateway API), StdoutEmitter (print to console).
\end{definition}


\section{Full Forge API}\label{sec:forge:api}

\begin{lstlisting}
// isls-forge/src/lib.rs

pub struct ForgeEngine {
    drill: DrillEngine,
    matrices: BTreeMap<String, Box<dyn Matrix>>,
    synthesizers: BTreeMap<String, Box<dyn Synthesizer>>,
    evaluators: Vec<Box<dyn Evaluator>>,
    emitters: Vec<Box<dyn Emitter>>,
    memory: PatternMemory,
    config: ForgeConfig,
}

pub struct ForgeConfig {
    pub pmhd: PmhdConfig,
    pub matrix: String,          // which matrix to use
    pub synth: String,           // which synthesizer
    pub emit: Vec<String>,       // which emitters
    pub validate: bool,          // run 8-gate validation
    pub output_dir: String,
}

pub struct ForgeResult {
    pub drill_result: DrillResult,
    pub artifacts: Vec<ForgeArtifact>,
    pub crystals: Vec<SemanticCrystal>,
    pub manifest: ExecutionManifest,
    pub pattern_count: usize,
}

pub struct ForgeArtifact {
    pub ir: ArtifactIR,
    pub synthesis: SynthesisOutput,
    pub evaluation: EvaluationResult,
    pub crystal_id: Hash256,
    pub emit_results: Vec<EmitResult>,
}

impl ForgeEngine {
    pub fn new(config: ForgeConfig, store: &IslandStore) -> Self;
    pub fn register_matrix(&mut self, m: Box<dyn Matrix>);
    pub fn register_synthesizer(&mut self, s: Box<dyn Synthesizer>);
    pub fn register_emitter(&mut self, e: Box<dyn Emitter>);
    pub fn forge(&mut self, spec: DecisionSpec) -> Result<ForgeResult>;
    pub fn forge_from_crystal(
        &mut self, crystal: &SemanticCrystal, domain: &str
    ) -> Result<ForgeResult>;
}
\end{lstlisting}


\section{CLI Integration}\label{sec:forge:cli}

\begin{lstlisting}
# Forge from a DecisionSpec
isls forge --spec intent.json --domain rust --output artifacts/

# Forge from an existing crystal
isls forge --crystal <crystal_id> --domain api --output artifacts/

# Forge with specific matrix
isls forge --spec intent.json --matrix workflow --output artifacts/

# List available matrices
isls forge --list-matrices

# Show pattern memory stats
isls forge --memory-stats
\end{lstlisting}


\section{Gateway Integration}\label{sec:forge:gateway}

\begin{lstlisting}[title={New Gateway Endpoints (added to C19)}]
POST /forge          Body: {spec, domain, config}
                     -> ForgeResult with crystal IDs and artifacts

POST /forge/crystal  Body: {crystal_id, domain}
                     -> ForgeResult from existing crystal

GET  /forge/matrices List available matrices
GET  /forge/memory   Pattern memory statistics
\end{lstlisting}


\section{Acceptance Tests}\label{sec:forge:tests}

\begin{longtable}{l l p{6.5cm}}
\toprule
\textbf{ID} & \textbf{Name} & \textbf{Procedure} \\
\midrule
\endfirsthead \midrule \endhead \bottomrule \endfoot
AT-F1 & End-to-end forge & Submit DecisionSpec for Rust module; verify
  file written and crystal committed. \\
AT-F2 & Forge determinism & Forge same spec twice; verify identical
  artifact content and crystal ID. \\
AT-F3 & Crystal validation & Forge an artifact; run validate formal;
  verify 8/8 checks pass. \\
AT-F4 & Forge from crystal & Discover crystal via analytical run; forge
  it into API spec; verify crystal committed. \\
AT-F5 & Matrix swap & Forge same spec with Rust matrix and API matrix;
  verify different output formats, both valid crystals. \\
AT-F6 & Constraint rejection & Submit spec with impossible constraint;
  verify evaluation fails and no crystal committed. \\
AT-F7 & Pattern memory persistence & Forge, restart, forge again;
  verify novelty score reflects the first run's patterns. \\
AT-F8 & Emitter file & Forge with FileEmitter; verify file exists at
  expected path. \\
AT-F9 & Emitter gateway & Start gateway; forge with GatewayEmitter;
  verify crystal queryable via GET /crystals. \\
AT-F10 & Custom matrix & Register custom matrix; forge with it;
  verify output uses custom interpretation. \\
\end{longtable}


%%% =====================================================================
\part{The Complete Loop}\label{part:loop}

\section{Discovery-Forge Cycle}\label{sec:loop}

With Phase~5, the complete \ISLS{} substrate supports:

\begin{enumerate}[nosep]
  \item \textbf{Observe}: Data flows in via adapters (C20) through the
        gateway (C19).
  \item \textbf{Discover}: The analytical engine (C1--C18) finds
        patterns and crystallizes them.
  \item \textbf{Forge}: A discovered crystal is fed into the forge
        (C21--C23) as a seed, combined with a DecisionSpec, and
        adversarially drilled into a monolith.
  \item \textbf{Synthesize}: The monolith becomes an ArtifactIR,
        interpreted through a domain matrix, synthesized into a
        concrete artifact.
  \item \textbf{Validate}: The artifact is packaged as a crystal and
        validated through the 8-gate cascade.
  \item \textbf{Emit}: The artifact is written to disk, sent to a
        webhook, or served via the gateway.
  \item \textbf{Deploy}: The artifact is used in the world.
  \item \textbf{Observe}: ISLS observes the effects of the deployed
        artifact.
  \item \textbf{Re-discover}: New patterns emerge from the changed
        world.
  \item \textbf{Re-forge}: The new patterns seed the next forge cycle.
\end{enumerate}

Every step in this cycle produces crystals. Every crystal is
content-addressed, evidence-bound, and replay-verifiable. The entire
history is append-only. The system accumulates structural intelligence
over time without manual configuration.


%%% =====================================================================
\part{Dependency Graph and Handoff}\label{part:handoff}

\section{Updated Dependency Graph}

\begin{lstlisting}[title={Complete ISLS Dependency Graph after Phase 5}]
isls-types        <-- (all crates)
    ^
isls-registry (C12)
    ^
isls-observe  <-- isls-persist <-- isls-topology (C16)
                       ^                ^
                       |          isls-extract
                       |                ^
                       |          isls-scale (C18)
                       |
               isls-consensus <-- isls-carrier
                       ^
               isls-archive <-- isls-morph
                       ^
               isls-manifest (C13)
                       ^
               isls-capsule (C14)
                       ^
               isls-scheduler (C15)
                       ^
               isls-engine
                       ^
               isls-store (C17)
                       ^
               isls-pmhd (C21)         <-- NEW: adversarial drill
                       ^
               isls-artifact-ir (C22)  <-- NEW: universal IR
                       ^
               isls-forge (C23)        <-- NEW: synthesis engine
                       ^
               isls-adapters (C20)
                       ^
               isls-gateway (C19)      <-- extended with /forge endpoints
                       ^
               isls-harness + isls-cli
\end{lstlisting}

\section{Test Count}

\begin{tabular}{l r r}
\toprule
\textbf{Phase} & \textbf{New Tests} & \textbf{Cumulative} \\
\midrule
Core (C1--C9) & 130 & 130 \\
Phase 1 (C10--C15) & 31 & 161 \\
Phase 2 (C16--C17) & 20 & 181 \\
Phase 3 (C18) & 15 & 196 \\
Phase 4 (C19--C20) & 23 & 219 \\
Phase 5 (C21--C23) & 22 & $\ge$ 241 \\
\bottomrule
\end{tabular}

\section{Verification Sequence}

\begin{lstlisting}
# Build
cargo build --workspace --release 2>&1 | grep -c warning  -> 0

# Test
cargo test --workspace -- --test-threads=1  -> >= 241 passed

# Forge from DecisionSpec
cat > /tmp/spec.json << 'EOF'
{
  "intent": "Create a REST API health-check module",
  "goals": {"coherence": 0.7, "robustness": 0.8},
  "constraints": ["must return JSON", "must include uptime"],
  "domain": "rust"
}
EOF
cargo run --bin isls --release -- forge --spec /tmp/spec.json \
  --domain rust --output /tmp/forge-out/

# Verify the forged artifact is a valid crystal
cargo run --bin isls --release -- validate formal
# -> Forge crystal passes 8/8 checks

# Forge from a discovered crystal
cargo run --bin isls --release -- ingest --adapter synthetic --scenario S-Basic
cargo run --bin isls --release -- run --mode shadow --ticks 100
CRYSTAL_ID=$(cargo run --bin isls --release -- crystal list | head -1)
cargo run --bin isls --release -- forge --crystal $CRYSTAL_ID \
  --domain api --output /tmp/forge-api/
\end{lstlisting}

\section{Completion Criteria}

Phase 5 is complete when:
\begin{enumerate}[nosep]
  \item PMHD drill produces monoliths deterministically.
  \item ArtifactIR builds deterministically from monoliths.
  \item Forge produces artifacts through all 4 built-in matrices.
  \item Every forged artifact is a valid crystal (8/8 checks).
  \item Pattern memory persists across runs.
  \item Forge from discovered crystal works end-to-end.
  \item Gateway /forge endpoints work.
  \item Total tests $\ge 241$, zero warnings.
  \item \textbf{23 Crates. One binary. Two directions. One substrate.}
\end{enumerate}


\vfill
\begin{center}
\rule{0.5\textwidth}{0.4pt}\\[0.5cm]
{\small Generated for \ISLS{} Extension Phase 5 v1.0.0.\\
The Eye discovers. The Hand forges. The Crystal proves.\\
Deterministic, append-only, replay-verified.}
\end{center}

\end{document}
